\begin{figure}[H]
\centering
\begin{tikzpicture}[
  >=Latex,
  every node/.style={font=\small},
  state/.style={draw,rounded corners=2pt,align=center,inner sep=5pt,text width=2.35cm,minimum height=.82cm},
  note/.style={draw,rounded corners=2pt,align=center,inner sep=5pt,text width=4.45cm,minimum height=.82cm},
  flow/.style={->,thick}
]
\node[font=\normalsize\bfseries] at (6.0,2.75) {仿射耦合：单向依赖带来可处理性};

\node[state] (x1) at (.8,1.20) {$x_1$\\上下文};
\node[state] (x2) at (.8,-.55) {$x_2$\\被变换};
\node[state] (st) at (4.0,.30) {缩放与平移\\由 $x_1$ 决定};
\node[state] (y1) at (7.35,1.20) {$y_1$\\保持不变};
\node[state] (y2) at (7.35,-.55) {$y_2$\\被变换};

\draw[flow] (x1) -- (st);
\draw[flow] (x1) -- (y1);
\draw[flow] (st) -- (y2);
\draw[flow] (x2) -- (y2);

\node[note] (tri) at (3.0,-2.45) {单向依赖\\$\Rightarrow$ 分块三角 Jacobian};
\node[note] (inv) at (8.7,-2.45) {保持不变的分支\\$\Rightarrow$ 可在逆变换中直接恢复};
\end{tikzpicture}
\caption{\textbf{仿射耦合中的非对称信息流。}
一部分变量原样通过，同时决定另一部分变量如何变换。这种单向依赖同时带来了三角 Jacobian 和闭式逆映射。}
\label{fig:affine-coupling-structure}
\end{figure}
